\documentclass[../../../include/open-logic-section]{subfiles}
% Part: incompleteness
% Chapter: incompleteness-provability
% Section: lob-thm

\begin{document}

\olfileid{inc}{inp}{lob}

\olsection{L\"ob定理}

理论~$\Th{T}$ 的Gödel语句是 $\lnot
\OProv[\Th{T}](y)$ 的一个不动点，即满足 \[
\Th{T} \Proves \lnot \OProv[\Th{T}](\gn{!G}) \liff !G.
\] 的语句~$!G$。
它是不可推导的，因为如果 $\Th{T} \Proves !G$，那么 (a) 由可证性条件~(1) 可得 $\Th{T} \Proves \OProv[\Th{T}](\gn{!G})$，
并且 (b) 将 $\Th{T} \Proves !G$ 与 $\Th{T} \Proves \lnot \OProv[\Th{T}](\gn{!G})
\liff !G$ 结合可得 $\Th{T} \Proves \lnot \OProv[\Th{T}](\gn{!G})$，
于是 $\Th{T}$ 就会是不一致的。
现在，很自然地会问：$\OProv[\Th{T}](y)$ 的不动点，即满足 \[
\Th{T} \Proves \OProv[\Th{T}](\gn{!H}) \liff !H.
\] 的语句~$!H$，地位又如何？
如果它是可推导的，那么由条件~(1) 可得 $\Th{T} \Proves \OProv[\Th{T}](\gn{!H})$，但对该等价式应用肯定前件律也能得出同样的结论。
因此，我们并不能得出 $\Th{T}$ 不一致，至少不能通过Gödel语句情形下的同样论证得出。
这当然并不表明 $\Th{T}$ \emph{确实}能推导出~$!H$。

对这个问题稍作推广，我们便能取得进展。
不动点等价式从左到右的方向，即 $\OProv[\Th{T}](\gn{!H}) \lif !H$，是被称为\emph{反射原理}的一般图式 $\OProv[\Th{T}](\gn{!A}) \lif !A$ 的一个实例。
它之所以这么叫，是因为它在某种意义上表达了 $\Th{T}$ 能够“反思”它所能推导的东西；大致来说，它断言：对任意~$!A$，“如果 $\Th{T}$ 能推导~$!A$，那么~$!A$ 就是真的。”
当然，这仅对可靠的理论成立，由此可知，理论通常不会推导出反射原理的每一个实例。
那么，一个（足够强且满足可证性条件的）理论能推导哪些实例呢？当然是所有那些 $!A$ 本身可推导的实例。仅此而已，正如下一个结果所示。

\begin{thm}\ollabel{thm:lob}
设 $\Th{T}$ 为扩张 $\Th{Q}$ 的可公理化理论，
并设 $\OProv[\Th{T}](y)$ 为满足
\olref[2in]{sec} 中条件 P1--P3 的公式。若 $\Th{T}$ 可推导 $\OProv[\Th{T}](\gn{!A}) \lif !A$，
则 $\Th{T}$ 事实上可推导 $!A$。
\end{thm}

换言之，如果 $\Th{T} \Proves/ !A$，那么 $\Th{T} \Proves/
\OProv[\Th{T}](\gn{!A}) \lif !A$。这一结果被称为 L\"ob 定理。

\begin{explain}
L\"ob 定理证明的启发式思路，在于一个关于圣诞老人存在的巧妙证明。
（如果你不喜欢这个结论，可以随意替换成任何你想要的结论。）证明如下：
\begin{enumerate}
\item 令 $X$ 为语句：“若 $X$ 为真，则圣诞老人存在。”
\item 假设 $X$ 为真。
\item 那么它所说的成立；即我们有：若 $X$ 为真，则圣诞老人存在。
\item 由于已假设 $X$ 为真，由 (2) 和 (3) 经肯定前件律可得：圣诞老人存在。
\item 我们成功从假设 (2) “$X$ 为真”导出了 (4) “圣诞老人存在”。由条件证明可知：“若 $X$ 为真，则圣诞老人存在。”
\item 但这恰是语句 $X$。因此，我们已表明 $X$ 为真。
\item 于是，由上述 (2)--(4) 的论证，圣诞老人存在。
\end{enumerate}
将这一想法形式化，把“为真”替换为“可推导”，把“圣诞老人存在”替换为 $!A$，
就得到了 L\"ob 定理的证明。技巧在于将不动点引理应用于公式~$\OProv[\Th{T}](y) \lif !A$。
该公式的不动点对应于前述梗概中的语句 $X$。
\end{explain}

\begin{proof}[证明 \olref{thm:lob}]
设 $!A$ 为使得 $\Th{T}$ 可推导 $\OProv[\Th{T}](\gn{!A}) \lif !A$ 的语句。
令 $!B(y)$ 为公式~$\OProv[\Th{T}](y)
\lif !A$，利用不动点引理可得语句 $!D$，使得 $\Th{T}$ 可推导 $!D \liff !B(\gn{!D})$。
那么下列每一项在 $\Th{T}$ 中都是可推导的：
\begin{align}
  & !D \liff (\OProv[\Th{T}](\gn{!D}) \lif !A) \ollabel{L-1}\\
  & \qquad \text{$!D$ 是~$!B(y)$ 的一个不动点}\notag \\
  & !D \lif (\OProv[\Th{T}](\gn{!D}) \lif !A) \ollabel{L-2}\\
  & \qquad\text{由 \olref{L-1}}\notag\\
  & \OProv[\Th{T}](\gn{!D \lif (\OProv[\Th{T}](\gn{!D}) \lif !A)}) \ollabel{L-3}\\
  & \qquad \text{由 \olref{L-2} 及条件 P1}\notag \\
  & \OProv[\Th{T}](\gn{!D}) \lif \OProv[\Th{T}](\gn{\OProv[\Th{T}](\gn{!D}) \lif !A})
  \ollabel{L-4}\\
  &\qquad \text{由 \olref{L-3} 使用条件 P2}\notag \\
  & \OProv[\Th{T}](\gn{!D}) \lif (\OProv[\Th{T}](\gn{\OProv[\Th{T}](\gn{!D})}) \lif \OProv[\Th{T}](\gn{!A})) \ollabel{L-5}\\
  &\qquad \text{由 \olref{L-4} 再次使用 P2} \notag\\
& \OProv[\Th{T}](\gn{!D}) \lif \OProv[\Th{T}](\gn{\OProv[\Th{T}](\gn{!D})}) \ollabel{L-6}\\
  & \qquad\text{由可证性条件 P3} \notag\\
  & \OProv[\Th{T}](\gn{!D}) \lif \OProv[\Th{T}](\gn{!A}) \ollabel{L-7} \\
  &\qquad\text{由 \olref{L-5} 与 \olref{L-6}}\notag\\
  & \OProv[\Th{T}](\gn{!A}) \lif !A \ollabel{L-8}\\
  &\qquad\text{由定理的假设} \notag\\
  & \OProv[\Th{T}](\gn{!D}) \lif !A \ollabel{L-9}\\
  &\qquad\text{由 \olref{L-7} 与 \olref{L-8}}\notag\\
  & (\OProv[\Th{T}](\gn{!D}) \lif !A) \lif !D \ollabel{L-10}\\
  & \qquad \text{由 \olref{L-1}}\notag \\
  & !D \ollabel{L-11}\\
  & \qquad\text{由 \olref{L-9} 与 \olref{L-10}}\notag \\
  & \OProv[\Th{T}](\gn{!D}) \ollabel{L-12}\\
  & \qquad\text{由 \olref{L-11} 及条件~P1}\notag \\
  & !A \qquad\qquad\text{由 \olref{L-8} 与 \olref{L-12}}\notag
\end{align}
\end{proof}

利用 L\"ob 定理，可以简短地证明第二不完全性定理（针对拥有满足条件 P1--P3 的可证性谓词的理论）：
如果 $\Th{T} \Proves
\OProv[\Th{T}](\gn{\lfalse}) \lif \lfalse$，那么 $\Th{T} \Proves \lfalse$。
如果 $\Th{T}$ 是一致的，则 $\Th{T} \Proves/ \lfalse$。
因此，$\Th{T}
\Proves/ \OProv[\Th{T}](\gn{\lfalse}) \lif \lfalse$，即 $\Th{T} \Proves/
\OCon[\Th{T}]$。
我们也可以应用该定理来证明 $\OProv[\Th{T}](x)$ 的不动点 $!H$ 是可推导的。
因为
\begin{align*}
  \Th{T} & \Proves \OProv[\Th{T}](\gn{!H}) \liff !H\\
  \intertext{特别地，有}
    \Th{T} & \Proves \OProv[\Th{T}](\gn{!H}) \lif !H
\end{align*}
于是由 L\"ob 定理，$\Th{T} \Proves !H$。

% Going in the other direction, for homework I
% may ask you to work through a short proof of L\"ob's theorem, using
% the second incompleteness theorem instead of the fixed-point lemma.

\begin{prob}
设 $\Th{T}$ 为可计算公理化的理论，
并令 $\OProv[\Th{T}]$ 为 $\Th{T}$ 的可证性谓词。
考虑以下四个陈述：
\begin{enumerate}
\item 若 $T \Proves !A$，则 $T \Proves \OProv[\Th{T}](\gn{!A})$。
\item $T \Proves !A \lif \OProv[\Th{T}](\gn{!A})$。
\item 若 $T \Proves \OProv[\Th{T}](\gn{!A})$，则 $T \Proves !A$。
\item $T \Proves \OProv[\Th{T}](\gn{!A}) \lif !A$
\end{enumerate}
每个陈述分别在什么条件下为真？
\end{prob}

\end{document}